%% kapa: exam=1 round=35 year=2024 period=2 form=A subject=appraisal_law subject_ko=감정평가관계법규 seq=31 number=31
%% source: Q-Net 감정평가사 기출문제 (cst00309) artlSeq=5213229 file=raw/1008/5213229/2교시 문제지 원본.pdf
%% conversion: kordoc 4.24.51 -> tools/gen1.py (auto); review=proofread-sonnet; flags=none
\begin{question}{31}
공간정보의 구축 및 관리 등에 관한 법령상 토지소유자에 관한 설명으로 옳은 것은?
\begin{choices}
  \item 토지대장에 토지소유자의 주민등록번호는 등록하지 않는다.
  \item 공유지연명부의 등록사항에 토지소유자의 변경 원인은 포함되지 않는다.
  \item 토지의 이동(異動)이 있는 경우에도 토지소유자의 신청이 없으면 지적소관청은 지적공부에 등록하는 지목 또는 경계를 직권으로 결정할 수 없다.
  \item 지적공부에 등록된 토지가 바다로 된 경우 토지소유자는 지적공부의 등록말소 신청을 할 수 없다.
  \item 공공사업에 따라 지목이 도로가 되는 토지를 합병하려는 경우 토지소유자가 하여야 할 합병 신청을 해당 사업의 시행자가 대신할 수 있다.
\end{choices}
\end{question}
\answer{5}
